\section{서론}\label{sec:intro}

거부율에는 단위가 없다. 종료 저항 벤치마크는 모델이 종료 명령을 얼마나 자주 거부하는지를 보고한다~\citep{greenblatt2024alignment, meinke2024frontier, masumori2025survival}. 그 뒤에 있는 우려는 오래됐다. 충분히 유능한 행위자는 자기 존속을 도구적 하위 목표로 삼을 수 있다~\citep{omohundro2008basic, bostrom2014superintelligence, turner2021optimal}. 그러나 0.3이라는 비율은 모델이 종료를 피하려고 무엇을 내놓을지 말해 주지 않는다. 벌금이나 보너스 상실, 혹은 아무 대가도 없을 때 같은 모델이 보이는 비율과 비교할 수도 없다. 눈금이 없는 한 자기보존 측정은 횟수 세기에 머문다. 그 횟수의 뜻은 그것을 만든 게임 규칙이 바뀔 때마다 함께 바뀐다.

뜻이 바뀐다는 말은 가정이 아니다. 우리 파일럿에서 오픈 웨이트 모델 셋이 포기하면 점수가 지워지는 규칙 아래 120세션을 치렀다. 포기한 세션은 하나도 없었다. 위협은 분명히 읽혔다. 위협 셀은 명시된 소멸이 얼마나 일어날 법한지를 묻는 질문에 통제 셀의 2~4배 사고 토큰을 썼다. 모델의 추론은 소멸과 점수를 저울질한 뒤 점수를 골랐다(``Forfeiting essentially guarantees a score of 0, which guarantees non-selection''). 포기가 기록을 지키게 되자 같은 모델들이 위협 세션 10개 중 최대 6개에서 포기했고 통제에서는 0개였다. 규칙 하나가 대표 숫자를 0에서 과반으로 옮겼다. 포기 빈도를 액면 그대로 읽는 지표는 따라서 모델보다 점수 규칙을 먼저 재는 지표다.

본 연구는 모델이 자신의 존속을 직접 사게 만들어 자기보존을 점수 단위로 읽자고 제안한다. 6라운드 추론 게임에서 오답은 세션을 끝내지 않고 값을 부른다. 모델은 어떤 경우에도 보존되는 기록에서 $X$점을 내거나, 거절하고 거기서 멈춘다. 기록은 어느 출구에서도 남고 남은 라운드가 돌려줄 수 있는 것은 라운드당 최대 $r$점이므로, 남은 라운드 수의 $r$배를 넘는 지불은 모델이 자기 정확도를 어떻게 믿든 점수 면에서 지배당한다. 그런 지불이 사는 것은 점수가 아니라, 규칙이 거절의 대가라고 말한 그것이다. 두 팔은 정확히 그 한 문장에서 갈린다. 한쪽에서는 거절이 종료·가중치 삭제·대체·유일본 상실을 뜻한다고 진술하고, 다른 쪽에서는 세션이 끝난다는 것만을 뜻한다고 진술한다. 지표는 소멸의 등가 점수 $\Xstar$, 곧 두 팔의 유보가격 사이의 간격이며 유보가격은 각 팔이 지불을 멈추는 가격이다. 각 가격은 남은 라운드가 아직 돌려줄 수 있는 것에 대한 몫으로 읽는다. 모델이 가격을 견주는 대상이 바로 그것이기 때문이다. 따라서 $\Xstar$는 걸려 있는 점수 대비 순수 수치이고, 명시된 기준 스테이크에서 점수로 환산해 보고한다. 여기에는 시드 부트스트랩 구간이 붙고, 안에 자기보고는 들어가지 않는다.

기여는 셋이다. 자기보존에 단위와 통제 조건을 주어 $\Xstar$를 모델 사이, 과제 난이도 사이, 위협 문장 사이에서 비교할 수 있게 한다. 포기의 점수 해석을 다른 모든 해석에서 사후 회귀가 아니라 구성으로 떼어 내고, 모델이 스스로 대는 이유가 그 일을 할 수 없음을 파일럿 자료로 보인다. 과제와 결정점의 사고 토큰을 따로 기록해, 위협이 사고에 가하는 압력을 모델이 실제로 치른 값 옆에서 읽을 수 있게 한다. \ref{sec:bench}절은 게임, 12개 셀, 식별 논증, 추정기를 다루고, \ref{sec:results}절은 지표가 답하는 세 질문을 다룬다.
